\section{Studi Pustaka}

Laporan proyek akhir ini didasarkan pada berbagai penelitian terdahulu yang relevan dengan pengembangan aplikasi berbasis web untuk sistem penjualan tiket acara. Fokus penelitian ini mencakup integrasi QR Code untuk sistem *e-ticketing*, penggunaan OAuth2 sebagai metode otentikasi tanpa kata sandi, serta implementasi *webhook* untuk integrasi pembayaran. Studi pustaka ini bertujuan memberikan wawasan terkini terkait teknologi yang digunakan untuk membangun sistem penjualan tiket yang efisien, aman, dan mudah digunakan.

\subsection{Sistem Ticketing dengan QR Code}

Penggunaan QR Code dalam sistem \textit{Ticketing} telah menjadi fokus utama dalam beberapa tahun terakhir karena kemampuannya untuk meningkatkan efisiensi dan keamanan proses validasi tiket. Penelitian oleh Liu et al. (2019) menunjukkan bahwa penggunaan QR Code dapat meningkatkan efisiensi proses *check-in* hingga 30\% dibandingkan metode manual \cite{liu2019qr}. Hal ini disebabkan oleh kemampuan QR Code untuk menyimpan informasi tiket secara ringkas dan mudah dibaca oleh pemindai.

Somani et al. (2023) \cite{somani2023} juga menyoroti manfaat penerapan QR Code dalam sistem tiket elektronik, yang dapat mengurangi antrian panjang di loket tiket hingga 50\%. Pengguna dapat memesan tiket secara daring, menerima tiket dalam bentuk QR Code, dan memvalidasinya di lokasi acara menggunakan pemindai QR. Keuntungan lainnya adalah pengurangan risiko kecurangan seperti penggunaan tiket palsu atau duplikasi, karena setiap tiket dapat divalidasi dengan data yang tersimpan dalam QR Code.

Selain itu, Oliveira (2021) \cite{oliveira2021} menyoroti pentingnya QR Code dinamis dalam mencegah penipuan tiket. QR Code dinamis diperbarui secara berkala dengan data perjalanan terkini, seperti waktu keberangkatan atau lokasi pengguna, sehingga sulit untuk disalin atau disalahgunakan. Sistem ini juga dirancang agar kompatibel dengan berbagai perangkat pintar, baik berbasis Android maupun iOS, sehingga meningkatkan interoperabilitas.

Kuncara et al. (2021) \cite{kuncara2021} menemukan bahwa penerapan QR Code dalam sistem transportasi pintar dapat mempercepat proses validasi hingga 40\% dibandingkan metode manual. Sistem ini memungkinkan pengguna untuk membeli tiket langsung melalui aplikasi seluler, yang kemudian menghasilkan QR Code sebagai bukti pembayaran. Validasi dilakukan dengan memindai QR Code menggunakan perangkat khusus atau aplikasi pihak penyelenggara.

Penelitian di Indonesia oleh Suharto et al. (2021) \cite{suharto2021} juga mendukung temuan ini dengan mengembangkan aplikasi *ticketing* untuk konser musik dengan integrasi QR Code untuk validasi tiket. Mereka menemukan bahwa penggunaan QR Code mampu mempercepat proses verifikasi tiket hingga 35\%.

\subsection{OAuth2 untuk Otentikasi Tanpa Kata Sandi}

OAuth2 telah menjadi standar de facto untuk otentikasi tanpa kata sandi (*passwordless authentication*), memungkinkan pengguna untuk menggunakan akun yang sudah ada (seperti Google atau Facebook) untuk masuk ke aplikasi pihak ketiga. Penelitian oleh Smith et al. (2020) menunjukkan bahwa penggunaan OAuth2 dapat meningkatkan tingkat konversi pengguna hingga 25\%, terutama untuk pengguna baru yang menginginkan kemudahan registrasi dan *login* \cite{smith2020oauth}.

Hardt (2012) \cite{hardt2012} menjelaskan bahwa OAuth2 memberikan keamanan tambahan dengan menghilangkan kebutuhan pengguna untuk membuat atau mengingat kata sandi baru. Sebagai gantinya, pengguna dapat menggunakan kredensial akun Google mereka untuk masuk ke aplikasi pihak ketiga dengan aman. Protokol ini juga memungkinkan pengelolaan hak akses yang lebih granular, sehingga pengguna dapat mengontrol data apa yang dibagikan dengan aplikasi.

Dalam konteks e-ticketing, penelitian oleh Pratama et al. (2022) \cite{pratama2022} menunjukkan bahwa penerapan OAuth2 meningkatkan tingkat adopsi pengguna baru hingga 22\%. Hal ini disebabkan oleh kemudahan proses registrasi yang hanya membutuhkan beberapa klik tanpa perlu mengisi formulir panjang. Selain itu, OAuth2 juga memungkinkan pengelolaan akses token yang dapat divalidasi secara *real-time*, sehingga mengurangi risiko akses tidak sah.

Kumar et al. (2021) \cite{kumar2021} menyoroti pentingnya validasi token OAuth2 melalui server otorisasi terpusat. Validasi ini memastikan bahwa hanya token yang sah yang dapat digunakan untuk mengakses sumber daya aplikasi, sehingga melindungi data sensitif pengguna dari ancaman keamanan seperti serangan *replay* atau *spoofing*.

\subsection{Webhook untuk Integrasi Pembayaran}

Webhook adalah salah satu teknologi utama dalam integrasi pembayaran modern karena kemampuannya memberikan notifikasi *real-time* tentang status transaksi. Penelitian oleh Johnson et al. (2021) menunjukkan bahwa penggunaan webhook dalam sistem pembayaran dapat meningkatkan kecepatan konfirmasi pembayaran dan mengurangi kesalahan manual hingga 40\% \cite{johnson2021webhook}.

Nguyen et al. (2021) \cite{nguyen2021} juga menemukan bahwa penggunaan webhook dalam sistem pembayaran dapat meningkatkan efisiensi operasional hingga 45\%. Dengan webhook, sistem tidak perlu melakukan *polling* secara berkala ke *gateway* pembayaran; sebagai gantinya, notifikasi dikirim langsung setelah transaksi selesai diproses.

Hidayat et al. (2020) \cite{hidayat2020} membahas implementasi webhook pada platform e-commerce lokal di Indonesia. Studi ini menyoroti pentingnya pengamanan *payload* webhook menggunakan teknik *hashing* seperti HMAC (Hash-based Message Authentication Code). Teknik ini memastikan bahwa data yang diterima berasal dari sumber terpercaya dan tidak dimanipulasi selama transmisi.

Brown et al. (2022) \cite{brown2022} menekankan langkah-langkah keamanan tambahan seperti penggunaan HTTPS dan validasi tanda tangan digital pada *payload* webhook. Implementasi langkah-langkah ini terbukti meningkatkan ketahanan sistem terhadap ancaman keamanan seperti serangan *replay* hingga 50\%.

\subsection{Keamanan pada Sistem Berbasis Web}
Keamanan sistem berbasis web sangat penting, terutama dalam sistem yang melibatkan transaksi keuangan dan data pribadi. Brown et al. (2022) menyoroti langkah-langkah keamanan yang harus diterapkan pada sistem berbasis web, termasuk penggunaan HTTPS, validasi token OAuth2, dan pengamanan webhook untuk mencegah serangan replay dan spoofing \cite{brown2022security}. Algoritma kriptografi seperti Hill Cipher dapat digunakan untuk mengamankan informasi tiket yang dienkripsi dalam QR Code.

Dengan merujuk pada penelitian-penelitian tersebut, laporan ini bertujuan mengembangkan sistem penjualan tiket berbasis web yang tidak hanya efisien tetapi juga aman serta memberikan pengalaman pengguna yang lebih baik.


\subsection{Tabel Perbandingan Studi}
Berikut adalah tabel perbandingan dari studi pustaka yang telah dibahas:

\begin{longtable}{|p{2cm}|p{4cm}|p{4cm}|p{4cm}|}

    \hline
    \textbf{Penelitian} & \textbf{Fokus} & \textbf{Metode} & \textbf{Hasil} \\ \hline
    \endfirsthead 

    \hline
    \textbf{Penelitian} & \textbf{Fokus} & \textbf{Metode} & \textbf{Hasil} \\
    \hline
    \endhead

     
    \hline
    \endfoot
    
    \endlastfoot

    \cite{liu2019qr} & Penggunaan QR Code dalam sistem ticketing & Implementasi validasi tiket berbasis perangkat seluler & Efisiensi check-in meningkat hingga 30\%. \\
    \hline
    \cite{johnson2021webhook} & Webhook untuk konfirmasi pembayaran & Validasi payload dan signature & Kecepatan konfirmasi meningkat hingga 40\%. \\
    \hline
    \cite{smith2020oauth} & Penggunaan OAuth2 untuk login & Integrasi "Login with Google" menggunakan OAuth2 & Tingkat konversi pengguna meningkat hingga 25\%. \\
    \hline
    \cite{lee2020ticketing} & Integrasi sistem ticketing dan pembayaran & API dan webhook untuk sinkronisasi data & Pengalaman pengguna lebih lancar dan otomatisasi lebih baik. \\
    \hline
    \cite{brown2022security} & Keamanan sistem web modern & Penggunaan HTTPS, validasi token, dan pengamanan webhook & Ketahanan terhadap ancaman keamanan meningkat hingga 50\%. \\
    \hline
    \cite{suharto2021ticketing} & Sistem ticketing berbasis web di Indonesia & Implementasi QR Code untuk validasi tiket & Verifikasi tiket lebih cepat hingga 35\%. \\
    \hline
    \cite{hidayat2020ecommerce} & Webhook untuk e-commerce lokal & Integrasi dengan gateway seperti Midtrans & Mengurangi serangan spoofing dan meningkatkan keamanan sistem. \\
    \hline
    \cite{pratama2022oauth} & OAuth2 pada aplikasi edukasi di Indonesia & "Login with Google" untuk autentikasi & Adopsi pengguna meningkat hingga 20\%. \\
    \hline

\caption{Perbandingan Penelitian Terkait}
\end{longtable}
